\begin{lemma}
\label{lem:p071-arbitrary-exact-star-polyhedral-insertion}
Let \(D\) be a compact standard core in an exact CLD star, with finite
exceptional set \(A\) and an arbitrary set \(B\setminus A\) of full
branches, and let \(Q\) be a nonempty compact convex polytope.  For
\(\phi\in C(D\times Q)\), let
\[
 C^{\mathrm{ex}}\phi
   =\sup\left\{v:v\leq\phi,\
       \substack{v\text{ is a continuous exact product-comparison}\\
                 \text{function on }D\times Q}\right\}.
\]
If \(u\leq\phi\) is any upper-semicontinuous exact product-comparison
function, then
\begin{equation}
                         u\leq C^{\mathrm{ex}}\phi.
 \label{eq:p071-arbitrary-star-envelope-dominates}
\end{equation}
If \(B\setminus A\) is countable, the continuous-cone envelope of
\Cref{lem:p071-exact-star-polyhedral-envelope} is consequently the largest
comparison minorant of \(\phi\), even when discontinuous comparison
minorants are allowed.
\end{lemma}

\begin{proof}
Fix \(x\in D\times Q\) and \(\delta>0\).  Continuity of \(\phi\) at the
node, uniformly over the compact set \(Q\), gives a finite set
\(F\subset B\setminus A\) such that
\begin{equation}
 |\phi_b(s,y)-\phi(\eta,y)|<\delta
 \quad
 \bigl(b\in B\setminus(A\cup F),\
       s\in[0,+\infty],\ y\in Q\bigr).
 \label{eq:p071-arbitrary-star-cofinite-obstacle}
\end{equation}
Indeed, otherwise a net of points on branches escaping every finite set,
together with a convergent subnet of their \(Q\)-coordinates, would
converge to a node point while violating continuity.  Enlarge \(F\) to
contain the full branch of \(x\), if \(x\) lies on one.

Restrict \(u\) and \(\phi\) to the finite exact substar consisting of all
exceptional branches and the branches in \(F\).
\Cref{lem:p071-finite-exact-star-polyhedral-insertion} gives a continuous
exact comparison function \(h\) there such that
\[
                              u\leq h\leq\phi.
\]
Extend \(h\) to the full core by setting every branch outside \(F\) equal
to its centre slice \(h_0(y)=h(\eta,y)\), and then subtract \(\delta\)
globally.  Call the result \(v\).  It is continuous at the node because
only finitely many branches are active, and it has the exact trace at every
full endpoint.  On an inactive branch,
\[
 v_b(s,y)=h_0(y)-\delta
       \leq\phi(\eta,y)-\delta
       <\phi_b(s,y)
\]
by \eqref{eq:p071-arbitrary-star-cofinite-obstacle}; on the finite substar
the inequality \(v\leq\phi\) is immediate.

We check the global junction inequality.  In a prescribed junction tuple,
split the sampled full branches into the active set \(S\subseteq F\) and
the inactive set \(T\).  Let \(\Lambda_0\) be the total reciprocal weight
of all exceptional and active samples, and let \(\Theta\) be that of the
inactive samples.  If \(\Lambda_0>0\), denote by \(y_*\) the barycentre of
the exceptional and active samples.  The finite-substar junction inequality
and convexity of \(h_0\) give
\[
 \begin{split}
  (\Lambda_0+\Theta)h_0(y_0)
  &\leq \Lambda_0h_0(y_*)
       +\sum_{b\in T}\frac{h_0(y_b)}{s_b}\\
  &\leq
       \sum_{i\in A\cup S}\frac{h_i(s_i,y_i)}{s_i}
       +\sum_{b\in T}\frac{h_0(y_b)}{s_b}.
 \end{split}
\]
If \(\Lambda_0=0\), the same conclusion is simply convexity of \(h_0\).
Global translation by \(-\delta\) preserves the inequality because the
reciprocal weights on both sides have the same sum.  Thus \(v\) is a
continuous exact product-comparison minorant of \(\phi\).

The chosen point belongs to the finite substar, so
\[
                         v(x)=h(x)-\delta\geq u(x)-\delta.
\]
It follows that
\(C^{\mathrm{ex}}\phi(x)\geq u(x)-\delta\).  Letting \(\delta\downarrow0\)
proves \eqref{eq:p071-arbitrary-star-envelope-dominates}.
\end{proof}
